% !TEX root = ../main.tex
\section{Synopsis}\label{sec:synopsis}

\textit{The Information Bottleneck (IB) theory of deep learning} is first 
introduced
in \cite{tishby2015}, building on the IB principle of \cite{tishby1999}, and
examined empirically in \cite{shwartzziv2017}. Here we will focus our attention
to the claims and results presented in \cite{shwartzziv2017}. They introduce a
promising model invariant methodology for each hidden layer $T_i$ in a DNN,
where each hidden layer $T_i$ is treated as one random variable, so that the
network forms the Markov chain $Y\to X\to T_1\to\dots\to T_k\to\hat{Y}$.
This measures what they in
\cite{shwartzziv2017} interpret as the compression and generalization 
capabilities
by respectively measuring the mutual information $I(X;T)$ and $I(Y;T)$. This
creates coordinates in the information plane \cite{tishby2015}.
The mutual information can't be measured directly, so they solve this by
binning the activations w.r.t. the $2^{12}=4096$
different input patterns into 30 equal interval bins. This is inevitable since 
$T$ is in principle a deterministic function of $X$ making the mutual 
information $I(X; T)$ infinite.

\textit{To argue the SGD Drift-diffusion phase}
they examine the behaviour of DNNs on said information plane as a function of
number of epochs. They claim
this provides evidence that training is split into two main phases; a short
initial fitting phase reducing training error, and one much longer compression
phase induced by stochastic relaxation or random diffusion. Thus owed to SGD
learning transitions into what they call the \textit{SGD drift-diffusion phase}
whose transition they claim is traceable both in the information plane, as
well as by the weight gradients mean norm and standard deviation.

\textit{The loss function of choice} for networks trained in 
\cite{shwartzziv2017} is
the \textit{cross-entropy loss function}. This choice is explained simply by
stating that it is conventional for stochastic optimization. They specifically
mention that they expect minimization of the Empirical Error Minimization (ERM)
to be - reasonably so - a result of the cross-entropy loss function. Seemingly
no more consideration as to a potential influence in compression is mentioned.

The cross-entropy loss or log-loss is defined as
\begin{equation}
	CE = -\sum_{i=1}^{N} q_i\log_2(p_i)
\end{equation}
